% GENERATED FILE. Edit canonical lessons, then run npm run generate:books.
% canonical-source-hash: fnv1a64:06d3342c13488db1
% canonical-lessons: UR-C174-makai, UR-C174-gandum, UR-C174-chane, UR-C174-rajma, UR-C174-anjir

\chapter{Maize, Wheat, Chickpeas, Kidney beans, Fig}
\label{ch:ur-maize-wheat-chickpeas-kidney-beans-fig}
\hlchaptermodality{174}

\begin{chapteropening}
\textbf{By the end of this chapter:} \emph{I can say maize, wheat, chickpeas, kidney beans and a fig.}

Complete the last lesson of chapter 174: I can say maize, wheat, chickpeas, kidney beans and a fig.
\end{chapteropening}

\section[\texorpdfstring{\ur{مکئی} (makaī)}{مکئی (makaī)}]{\ur{مکئی} (makaī) — maize}

\label{lesson:UR-C174-makai}



\begin{quote}
Before the new one: say the Urdu for a tree trunk, then the Urdu for mud.
\end{quote}

\subsection*{You'll want to know: \ur{مکئی}}
\textbf{\ur{مکئی}} — \emph{makaī} — \textquotedblleft{}maize\textquotedblright{}.

\begin{cousinweb}
CDIAL lists \textbf{\ur{مکئی}} (\emph{makaī}), as Hindi \emph{makāī}, \textquotedblleft{}maize\textquotedblright{}, under Sanskrit \emph{markaka}, a bird-name known only from the old word-lists, without explaining the link (CDIAL 9879).
\end{cousinweb}

\begin{grammarlens}[title={What you've built}]
One more word for this chapter.
\end{grammarlens}

\subsection*{Your turn}
\begin{itemize}
\raggedright
  \item \emph{Say it:} \emph{makaī}
  \item \emph{Say it:} \emph{makaī}, once more
  \item \emph{Say it:} say \emph{makaī}
  \item \emph{Recall:} say the Urdu for a tree trunk, then the Urdu for mud, then say \emph{makaī} again
\end{itemize}

\begin{tcolorbox}[breakable,colback=teal!4,colframe=teal!35!black,title={Before you move on}]
What does \ur{مکئی} mean? (\textquotedblleft{}Maize\textquotedblright{}.) Say it once more.
\end{tcolorbox}
\section[\texorpdfstring{\ur{گندم} (gandum)}{گندم (gandum)}]{\ur{گندم} (gandum) — wheat}

\label{lesson:UR-C174-gandum}



\begin{quote}
Before the new one: say the Urdu for mud, then the Urdu for maize.
\end{quote}

\subsection*{You'll want to know: \ur{گندم}}
\textbf{\ur{گندم}} — \emph{gandum} — \textquotedblleft{}wheat\textquotedblright{}.

\begin{cousinweb}
\textbf{\ur{گندم}} (\emph{gandum}) is a loan, usually traced to Persian \emph{gandum}, \textquotedblleft{}wheat\textquotedblright{}.
\end{cousinweb}

\begin{grammarlens}[title={What you've built}]
One more word for this chapter.
\end{grammarlens}

\subsection*{Your turn}
\begin{itemize}
\raggedright
  \item \emph{Say it:} \emph{gandum}
  \item \emph{Say it:} \emph{gandum}, once more
  \item \emph{Say it:} say \emph{gandum}
  \item \emph{Recall:} say the Urdu for mud, then the Urdu for maize, then say \emph{gandum} again
\end{itemize}

\begin{tcolorbox}[breakable,colback=teal!4,colframe=teal!35!black,title={Before you move on}]
What does \ur{گندم} mean? (\textquotedblleft{}Wheat\textquotedblright{}.) Say it once more.
\end{tcolorbox}
\section[\texorpdfstring{\ur{چنے} (chane)}{چنے (chane)}]{\ur{چنے} (chane) — chickpeas}

\label{lesson:UR-C174-chane}



\begin{quote}
Before the new one: say the Urdu for maize, then the Urdu for wheat.
\end{quote}

\subsection*{You'll want to know: \ur{چنے}}
\textbf{\ur{چنے}} — \emph{chane} — \textquotedblleft{}chickpeas\textquotedblright{}.

\begin{cousinweb}
\textbf{\ur{چنے}} (\emph{chane}) is a tadbhava, a word inherited from Sanskrit and reshaped by sound change. It comes from \emph{caṇa}, \textquotedblleft{}chickpea\textquotedblright{} (CDIAL 4579). \emph{chane} is the plural of \emph{chanā}.
\end{cousinweb}

\begin{grammarlens}[title={What you've built}]
One more word for this chapter.
\end{grammarlens}

\subsection*{Your turn}
\begin{itemize}
\raggedright
  \item \emph{Say it:} \emph{chane}
  \item \emph{Say it:} \emph{chane}, once more
  \item \emph{Say it:} say \emph{chane}
  \item \emph{Recall:} say the Urdu for maize, then the Urdu for wheat, then say \emph{chane} again
\end{itemize}

\begin{tcolorbox}[breakable,colback=teal!4,colframe=teal!35!black,title={Before you move on}]
What does \ur{چنے} mean? (\textquotedblleft{}Chickpeas\textquotedblright{}.) Say it once more.
\end{tcolorbox}
\section[\texorpdfstring{\ur{راجما} (rājmā)}{راجما (rājmā)}]{\ur{راجما} (rājmā) — kidney beans}

\label{lesson:UR-C174-rajma}



\begin{quote}
Before the new one: say the Urdu for wheat, then the Urdu for chickpeas.
\end{quote}

\subsection*{You'll want to know: \ur{راجما}}
\textbf{\ur{راجما}} — \emph{rājmā} — \textquotedblleft{}kidney beans\textquotedblright{}.

\begin{cousinweb}
\textbf{\ur{راجما}} (\emph{rājmā}) is often linked with Sanskrit \emph{rājamāṣa}, \textquotedblleft{}a kind of bean\textquotedblright{}, built from \emph{rājan}, \textquotedblleft{}king\textquotedblright{}, and \emph{māṣa}, \textquotedblleft{}bean\textquotedblright{}. CDIAL records that word, but lists no Hindi form under it (CDIAL 10684).
\end{cousinweb}

\begin{grammarlens}[title={What you've built}]
One more word for this chapter.
\end{grammarlens}

\subsection*{Your turn}
\begin{itemize}
\raggedright
  \item \emph{Say it:} \emph{rājmā}
  \item \emph{Say it:} \emph{rājmā}, once more
  \item \emph{Say it:} say \emph{rājmā}
  \item \emph{Recall:} say the Urdu for wheat, then the Urdu for chickpeas, then say \emph{rājmā} again
\end{itemize}

\begin{tcolorbox}[breakable,colback=teal!4,colframe=teal!35!black,title={Before you move on}]
What does \ur{راجما} mean? (\textquotedblleft{}Kidney beans\textquotedblright{}.) Say it once more.
\end{tcolorbox}
\section[\texorpdfstring{\ur{انجیر} (anjīr)}{انجیر (anjīr)}]{\ur{انجیر} (anjīr) — a fig}

\label{lesson:UR-C174-anjir}



\begin{quote}
Before the new one: say the Urdu for chickpeas, then the Urdu for kidney beans.
\end{quote}

\subsection*{You'll want to know: \ur{انجیر}}
\textbf{\ur{انجیر}} — \emph{anjīr} — \textquotedblleft{}a fig\textquotedblright{}.

\begin{cousinweb}
\textbf{\ur{انجیر}} (\emph{anjīr}) is usually traced to Persian \emph{anjīr}, \textquotedblleft{}fig\textquotedblright{}. CDIAL lists Hindi \emph{añjīr} under Sanskrit \emph{añjīra}, \textquotedblleft{}fig tree\textquotedblright{}, and marks that word as itself taken from Persian (CDIAL 174).
\end{cousinweb}

\begin{grammarlens}[title={What you've built}]
One more word for this chapter.
\end{grammarlens}

\subsection*{Your turn}
\begin{itemize}
\raggedright
  \item \emph{Say it:} \emph{anjīr}
  \item \emph{Say it:} \emph{anjīr}, once more
  \item \emph{Say it:} say \emph{anjīr}
  \item \emph{Recall:} say the Urdu for chickpeas, then the Urdu for kidney beans, then say \emph{anjīr} again
\end{itemize}

\begin{tcolorbox}[breakable,colback=teal!4,colframe=teal!35!black,title={Before you move on}]
What does \ur{انجیر} mean? (\textquotedblleft{}A fig\textquotedblright{}.) Say it once more.
\end{tcolorbox}
